%%%%%%%%%%  scan idx 49  |  folio 38 (continued)  %%%%%%%%%%

\underline{Definición 1.7. Para cada objeto $R$ de
$\underline C_\Lambda$, pongamos}
\[
 \overline\Omega_R=M/(m_\Lambda R+M^2),
\]
\underline{donde} $\underline{m}_\Lambda$ \underline{y $M$ son los ideales maximales de
$\Lambda$ y $R$, respectivamente. Este es un espacio vectorial de
dimensión finita sobre $k$, y $R\longmapsto\overline\Omega_R$ define
un funtor}
\[
 \overline\Omega:\underline C_\Lambda
   \longrightarrow\underline{\underline V},
\]
\underline{donde $\underline{\underline V}$ es la categoría de los
espacios vectoriales de dimensión finita sobre $k$. Si $\underline A$
es una categoría cofibrada sobre $\underline C_\Lambda$, el funtor
compuesto $\underline A\xrightarrow P\underline C_\Lambda
\longrightarrow\underline{\underline V}$ se denota también por
$\overline\Omega$, por abuso de notación.}

\underline{Proposición 1.8. Sea $\underline A$ un grupoide cofibrado
semihomogéneo sobre $\underline C_\Lambda$. Dado un objeto $a$ de
$\underline A$ que es versal para las pequeñas prolongaciones de $e$,
existe una flecha $a'\longrightarrow a$ tal que:}

%%%%%%%%%%  scan idx 50  |  folio 39  %%%%%%%%%%

(i) \underline{$a'$ es versal para las pequeñas prolongaciones de $a$;}

(ii) \underline{$P(a')\longrightarrow P(a)$ es sobreyectivo, y su
núcleo está anulado por el ideal maximal de $P(a)$;}

(iii) \underline{$\overline\Omega_{a'}\longrightarrow
\overline\Omega_a$ es un isomorfismo.}

\underline{Demostración.} Pongamos $R=P(a)$ y consideremos una sucesión
exacta
\[
 0\longrightarrow J\longrightarrow S\longrightarrow R
 \longrightarrow0,
\]
donde $S$ es un anillo de series formales sobre $\Lambda$ en un número
finito de variables y $J\subset M^2$, siendo $M$ el ideal maximal de
$S$. Sea $\Sigma$ el conjunto de los ideales $J_1$ de $S$ tales que

(i) $J\supset J_1\supset MJ$, y

(ii) existe un $\pi_1$-morfismo $a_1\longrightarrow a$ en
$\underline A$, donde $\pi_1:S/J_1\longrightarrow R$ es el morfismo
canónico.

La propiedad 1.2(1) de $\underline A$, junto con el hecho de que $J/MJ$
es un espacio vectorial de dimensión finita sobre $k$, implica que
$\Sigma$ tiene un ideal mínimo. Denotémoslo por $J'$ y elijamos un
$\pi'$-morfismo $a'\xrightarrow g a$, donde
$\pi':S/J'\longrightarrow R$ es el morfismo canónico. Como
$J'\subset J\subset M^2$, el morfismo
$\overline\Omega_{a'}\longrightarrow\overline\Omega_a$ es un
isomorfismo.

Afirmamos que $a'\xrightarrow g a$ es versal para las pequeñas
prolongaciones de $a$. En efecto, sea $a_1\longrightarrow a$ una
pequeña prolongación y pongamos $R_1=P(a_1)$. Obtenemos un diagrama
conmutativo exacto
\[
\begin{tikzcd}[column sep=small]
0 \arrow[r] & J/MJ \arrow[r] \arrow[d,"h"]
 & S/MJ \arrow[r] \arrow[d,"h_1"]
 & R \arrow[r] \arrow[d,equal] & 0 \\
0 \arrow[r] & k \arrow[r] & R_1 \arrow[r] & R \arrow[r] & 0.
\end{tikzcd}
\]
Si $h=0$, entonces $\pi_1:R_1\longrightarrow R$ admite una sección,
por lo que $h_1$ induce un morfismo $f:S/J'\longrightarrow R_1$. Si
$h\ne0$, entonces $h_1:S/MJ\longrightarrow R_1$ es sobreyectivo y,
por tanto, por la definición de $J'$, induce de nuevo un morfismo
$f:S/J'\longrightarrow R_1$. En ambos casos obtenemos un diagrama
conmutativo
\[
\begin{tikzcd}[column sep=large]
S/J' \sisi{\arrow[rr,"f",no head]}{\arrow[rr,"f"]} \arrow[dr,"\pi'"']
 & & R_1 \arrow[dl,"\pi_1"] \\
& R &
\end{tikzcd}
\]
Como $a$ es versal para las pequeñas prolongaciones de $e$, el
Corolario 1.6 proporciona un diagrama conmutativo

%%%%%%%%%%  scan idx 51  |  folio 40  %%%%%%%%%%

\[
\begin{tikzcd}[column sep=large]
a' \arrow[rr,"\rho"] \arrow[dr] & & a_1 \arrow[dl] \\
& a &
\end{tikzcd}
\]
como se quería.
